\makeabstract
{
Is Lake Michigan a Metric (Dog-)Space?
}
{
Oscar Hernandez (1), Grace Stoner (2), Matt Glomski, PhD (2), Izzy Petersen (3)
}
{
\textsuperscript{1}Department of Mathematics, Amherst College, Amherst, MA, USA\\
\textsuperscript{2}Department of Mathematics, Marist University, Poughkeepsie, NY, USA\\
\textsuperscript{3}Department of Mathematics, Union College, Schenectady, NY, USA
}
{
In 2003, the mathematics community was introduced to a Welsh Corgi named Elvis in the paper "Do dogs know calculus?" In that article, author Timothy J. Pennings describes a game of fetch he played with his dog at Lake Michigan. With Elvis standing next to him at the shoreline, the author would throw a tennis ball into the water. The shortest route is a straight swim, while any route that involves running can be covered at a greater average speed than pure swimming. The problem for Elvis — and for generations of first-semester calculus students — is to find the precise path that minimizes the total travel time to the ball. Our research aimed to generalize the problem to include more paths.
}
{
As a problem rooted in distance, we sought to describe the problem in terms of a distance function, otherwise known as a metric. After describing a metric, we aimed to use that metric to numerically approximate the region Elvis could cover. The region that describes where Elvis can reach in x amount of time from point y is called the open ball of our metric beach-water space. 
}
{
The results of our research show that Elvis’ path, be it water-to-water, land-to-water, or land-to-land, result in an extraordinarily shaped open ball. We learned that the shape was significantly more complicated than previously thought, requiring an non-elemental curve to describe part of the open ball.
}
{
In conclusion, our research investigated the optimal paths Elvis could take from various starting positions. Our work resulted in an open ball with an unconventional shape. Further steps would include determining the precise curve that outlines the open ball. Additionally, another point of investigation would be to create a more accurate metric that models the continuous velocity change at the shoreline.
}

\newpage
